\subsection{Private operator commitments and the versioned instance}
\label{sec:private-operator-instance}

The private operator opening consists of the canonical coupling packet and
the separate 32-byte nonce \(r\). In the fixed scalar Euclidean profile,
this packet determines the cellular maps, Hodge operator, integral structure
and intrinsic harmonic lattice. Define
\[
 \operatorname{Com}(B;r)=
 H\bigl(\mathsf{dom}_{\mathrm{op}}\Vert
   \operatorname{frame}(\operatorname{enc}(B))\Vert
   \operatorname{frame}(r)\bigr).
\]
The nonce is distinct from the source credential salt. Canonical decoding
enforces the operator domain, original matrix codec, fixed nonce width and
full consumption. Distinct accepted matrices with equal commitments give
an actual hash collision between these complete preimages, allowing their
nonces to differ.

For the private carrier \(w=(S,B,z,a,u,\rho,r)\) and fixed descriptor
\(p=(m,n,K,\ell)\), define its public instance by the three names
\begin{equation}
 x_w=\bigl(H(\operatorname{credbytes}(\rho,S)),\,
          \operatorname{Com}(B;r),\,H(\beta)\bigr),
 \qquad \beta=E_B(z,a,u).
 \label{eq:coupled-committed-instance}
\end{equation}
Only the instance and ciphertext bytes are released; neither matrix nor
opening is serialized in the public instance. The common declared semantic
class is fixed for these dimensions. The versioned private-witness,
program and instance formats use the respective domains
\texttt{HODGE\_NATIVE\_WITNESS\_V2},
\texttt{HODGE\_NATIVE\_PROGRAM\_V2} and
\texttt{HODGE\_NATIVE\_INSTANCE\_V2}. They retain the canonical matrix,
integer and salt encodings and reject the earlier raw-operator-hash formats.
The finite Boolean relation is
\[
\begin{aligned}
 R_{p,H}(x,\eta)=\mathsf{true}
 \quad\Longleftrightarrow\quad \exists w:\;&
 \eta=\operatorname{encodeWitness}(w),\\
 &w\text{ satisfies the bounds for }p,\quad x=x_w.
\end{aligned}
\]
Canonical program decoding fixes \(m,n,K,\ell\), and the checker recomputes
all three names from the decoded witness. This gives both honest-witness
acceptance and reflection of acceptance into the stated relation.

\begin{proposition}[Committed-instance recovery or a hash collision]
\label{prop:coupled-committed-recovery}
Suppose an accepted witness \(w\) names an original target \(B_0\), with
operator nonce \(r_0\), and the original ciphertext
\(E_{B_0}(z_0,a_0,u_0)\). Then its message is \(z_0\), or there are distinct
canonical operator preimages with equal hashes, or distinct canonical
ciphertext byte strings with equal hashes.
\end{proposition}

\begin{proof}
The second name gives equality of the two target commitments. Unless their
complete preimages exhibit a hash collision, canonical injectivity gives
\(B=B_0\). The third name similarly gives equality of the ciphertext bytes
unless they exhibit a byte-hash collision. In the remaining case,
Proposition~\ref{prop:coupled-byte-recovery} applied with this same private
matrix recovers both messages from the same bytes, so \(z=z_0\).
\end{proof}

For the signed family, let \(\varepsilon(0)=-1\) and \(\varepsilon(1)=1\), applied
coordinatewise to Boolean vectors, and put
\[
 B_k=A\operatorname{diag}(\varepsilon(k)),\qquad
 z=d\varepsilon(\mu),\qquad u=B_k\varepsilon(\mu).
\]
Sample the Boolean message \(\mu\in\{0,1\}^n\), key \(k\in\{0,1\}^n\)
and nonce independently and uniformly. Fix the source opening \((\rho,S)\)
and exact opening \(a\) independently of these samples. Its credential is
therefore fixed in the following fixed-hash experiment. The released view
is the pair
\[
 V=\bigl(\operatorname{encodeInstance}(x_w),E_{B_k}(z,a,u)\bigr).
\]
If the bounded checker is used, select \(K\) large enough for this finite
family and the fixed source opening.

Fix a reference matrix \(B_*\) of the same dimensions. For a deterministic
complete-message decoder \(D\), the chosen-operator test \(T_D\) samples
its own \(\mu,k\), receives either \(\operatorname{Com}(B_k;r)\) or
\(\operatorname{Com}(B_*;r)\), constructs the same ciphertext and public
view, runs \(D\), and tests whether the result is \(\mu\). Let \(p_1,p_0\)
be its real and reference acceptance probabilities, respectively, and put
\(\operatorname{Adv}_{\mathrm{op}}(T_D)=|p_1-p_0|\).

\begin{theorem}[Public-instance recovery reduces to operator hiding]
\label{thm:coupled-instance-recovery}
For \(A^{\mathsf T}A=dI\) and every such decoder \(D\),
\[
 \Pr[D(V)=\mu]\leq 2^{-n}+\operatorname{Adv}_{\mathrm{op}}(T_D).
\]
For positive \(d\), a partial integer-coefficient extractor has the same
bound using its threshold-totalized bit decoder in \(T_D\).
\end{theorem}

\begin{proof}
In the reference world, the operator root depends only on the independent
nonce and fixed matrix. The ciphertext law, with that nonce as an
independent tape, gives complete-message recovery probability exactly
\(2^{-n}\). The real world's public names and ciphertext are exactly those
computed from the sampled private witness. The two recovery events are the
acceptance events of \(T_D\), so their difference is bounded by its
advantage. For a partial coefficient extractor, replace failure by a fixed
bit vector and an integer output by its coordinatewise positive-sign bits.
When \(d>0\), every correct coefficient output gives the correct Boolean
message. Its success event is therefore contained in the totalized decoder's
success event.
\end{proof}

This is an algebraic transfer for a fixed byte hash, with the advantage of
an explicit chosen-operator test. The next experiment instead samples one
shared random oracle and bounds the actual recovery event directly.
